% Einheit 3 von 5 – Prisma (bank/koerper/e3.jsonl, zusammenbau v0.1)
%% TODO Zweigzeile Teil 1: was hier gelernt wird (Satz in Schülersprache aus dem Einheitstitel) – Einheit 3
\einheitenkopf[e3]{Einheit 3 von 5 · Prisma}
\zweigzeile{ab Kl. 8 (am Gymnasium ab Kl. 7) · P10}
\setcounter{aufgabe}{19}

%% TODO Titel: Ich-kann-Satz fehlt in der Bank (Platzhalter: Kettenname) – Nr. 20
%% TODO Anweisung: ein Satz über der ersten Teilaufgabe fehlt in der Bank – Nr. 20
\begin{aufgabe}{Prisma}
\swa{Ein Zelt hat die Form dieses Prismas. Färbe die Grundfläche und zeichne die Höhe des Prismas farbig nach.}{\prismadreieck{4}{2.5}{2}{}{}{}}
%% TODO Darstellung für schwach fehlt in der Bank (koerper-e3-k1-s1-v1; Abschnitt „Für schwache Schüler“)
\swz{Prisma: Grundfläche $18$ cm², Höhe $7$ cm – Volumen?}{}
%% TODO Darstellung für schwach fehlt in der Bank (koerper-e3-k1-s1-v2; Abschnitt „Für schwache Schüler“)
\swz{Prisma: Grundfläche $25$ cm², Höhe $6$ cm – Volumen?}{}
%% TODO Darstellung für schwach fehlt in der Bank (koerper-e3-k1-s1-v3; Abschnitt „Für schwache Schüler“)
\swz{Prisma: Grundfläche $32$ dm², Höhe $9$ dm – Volumen?}{}
%% TODO Darstellung für schwach fehlt in der Bank (koerper-e3-k1-s1-v4; Abschnitt „Für schwache Schüler“)
\swz{Ein Modell eines Deichs hat die Form eines Prismas. Die Grundfläche ist ein gleichschenkliges Trapez mit dem Flächeninhalt $2\,340$ cm², die Höhe des Prismas ist $85$ cm. Berechne das Volumen \hfill (P10 2014 OS).}{}
%% TODO Darstellung für schwach fehlt in der Bank (koerper-e3-k1-s1-v5; Abschnitt „Für schwache Schüler“)
\swz{Ein Blumenkasten hat die Form eines Prismas. Die Grundfläche ist ein gleichschenkliges Trapez mit dem Flächeninhalt $418$ cm², die Höhe des Prismas ist $72$ cm. Berechne das Volumen \hfill (P10 2014 OS).}{}
\swfrage{Was bleibt gleich, was ändert sich?}
%% TODO Darstellung für schwach fehlt in der Bank (koerper-e3-k1-s2-v1; Abschnitt „Für schwache Schüler“)
\swz{Prisma mit rechteckiger Grundfläche $7$ cm × $3$ cm, Höhe $5$ cm – Volumen?}{}
%% TODO Darstellung für schwach fehlt in der Bank (koerper-e3-k1-s3-v1; Abschnitt „Für schwache Schüler“)
\swz{Dreiecksprisma: Dreieck mit Grundseite $8$ cm und Höhe $5$ cm; Höhe des Prismas $9$ cm – Volumen?}{}
%% TODO Darstellung für schwach fehlt in der Bank (koerper-e3-k1-s4-v1; Abschnitt „Für schwache Schüler“)
\swz{Prisma mit Trapez als Grundfläche: parallele Seiten $7$ cm und $3$ cm, Abstand $4$ cm; Höhe des Prismas $6$ cm – Volumen?}{}
\swz{Ein Zelt liegt wie im Bild: vorn ein Dreieck, $2{,}4$ m breit und $1{,}6$ m hoch; das Zelt ist $3$ m lang. Wie viel Luft ist im Zelt?}{\prismadreieck{4}{2.5}{2}{$2{,}4$ m}{$1{,}6$ m}{$3$ m}}
%% TODO Darstellung für schwach fehlt in der Bank (koerper-e3-k1-s6-v1; Abschnitt „Für schwache Schüler“)
\swz{Die Grundfläche eines Prismas ist ein Dreieck mit den Seiten $4$, $7$ und $8$ cm, das Prisma ist $5$ cm hoch – Mantel?}{}
%% TODO Darstellung für schwach fehlt in der Bank (koerper-e3-k1-s7-v1; Abschnitt „Für schwache Schüler“)
\swz{Dreiecksprisma: Das Dreieck hat einen rechten Winkel zwischen den Seiten $3$ cm und $4$ cm, die lange Seite ist $5$ cm; das Prisma ist $7$ cm hoch – Oberfläche?}{}
\swa{Vom Netz eines Dreiecksprismas sind drei Rechtecke nebeneinander gezeichnet: $4$ cm × $7$ cm, $5$ cm × $7$ cm und $6$ cm × $7$ cm. Zeichne das Netz ab und ergänze, was fehlt.}{\rechenplatz{5}}
%% TODO Darstellung für schwach fehlt in der Bank (koerper-e3-k1-s9-v1; Abschnitt „Für schwache Schüler“)
\swz{Prisma: Volumen $210$ cm³, Grundfläche $35$ cm² – Höhe?}{}
\swz[3]{Ein Brunnentrog hat die Form eines Prismas; Grund- und Deckfläche sind gleichseitige Dreiecke mit der Seite $1{,}80$ m und der Höhe $1{,}56$ m. Das Prisma ist $0{,}75$ m hoch. Berechne den Flächeninhalt der Deckfläche und weise nach, dass das Volumen etwa $1{,}1$ m³ beträgt \hfill (P10 2019 OS).}{}
\end{aufgabe}

%% TODO Titel: Ich-kann-Satz fehlt in der Bank (Platzhalter: Kettenname) – Nr. 21
%% TODO Anweisung: ein Satz über der ersten Teilaufgabe fehlt in der Bank – Nr. 21
\begin{aufgabe}{Prisma – Fehler finden, Begründen, Anwendung, Darstellungswechsel}
\swz[3]{Dreiecksprisma: Dreieck mit Grundseite $7$ cm und Höhe $6$ cm, das Prisma ist $5$ cm hoch. Emil rechnet: \rechnung{7 \cdot 6 &= 42 \text{ cm}^2 \\ 42 \cdot 5 &= 210 \text{ cm}^3} Wo ist der Fehler?}{}
\swfrage{Warum rechnet man beim Prisma Grundfläche mal Höhe? Denk an einen Stapel.}
\swz[3]{Ein Dachboden hat die Form eines liegenden Dreiecksprismas: $9$ m breit, $3{,}5$ m hoch und $11$ m lang. Ein Heizgerät reicht für Räume bis $180$ m³. Reicht es für den ausgebauten Dachboden?}{}
\swa{Das Dreieck ist gleichschenklig, die Schenkel sind je $5$ cm lang. Skizziere das Netz des Prismas mit allen Maßen.}{\prismadreieck{4}{1.5}{2.5}{$8$ cm}{$3$ cm}{$7$ cm} \rechenplatz[halb]{4}}
\end{aufgabe}

\uebersichtskasten{Volumen: Grundfläche mal Höhe – die Grundfläche ist die Fläche, die vorn und hinten gleich ist (auch wenn das Prisma liegt). \\ Dreieck g = 6 cm, h = 4 cm, Prismenhöhe 10 cm: G = 6 · 4 : 2 = 12 cm² \quad V = 12 · 10 = 120 cm³ \\ Mantel: alle Rechtecke zusammen – Umfang der Grundfläche mal Höhe. \quad u = 6 + 5 + 5 = 16 cm: M = 16 · 10 = 160 cm² \\ Oberfläche: zweimal Grundfläche plus Mantel. \quad O = 2 · 12 + 160 = 184 cm²}
